\scanpage{143}
% archive sheet number: 142
\mspage{1.}

\S\,1 \uline{Notation and terminology}: Let $A$ be a Dedekind ring and let $\overline A$ be the maximal unramified extension of $A$. Let $G_A$ be the Galois group of $\overline A$ over $A$ (fundamental group of $\operatorname{Spec} A$). Let $M$ be a \uline{finite} $G_A$-module whose order, $m=\operatorname{Card} M$, is prime to all residue class characteristics of $A$, i.e. such that $\operatorname{Card} M$ is invertible in $A$. Suppose that $G_A$ operates continuously on $M$, i.e. that some open subgroup operates trivially. If $G_A$ is \uline{infinite}, we put
\[
H^r(A,M)=H^r(G_A,M)=\lim_{\underset B\longrightarrow} H^r(G_{B/A},M),
\]
the limit being taken over all \uline{finite} unramified extensions $B/A$ with respect to the inflation maps. If $G_A$ is \uline{finite}, (which for us will occur only if $A$ is the field of real or complex numbers) we put
\[
H^r(A,M)=\underline H^r(G_A,M),
\]
where the underline denotes the ``complete'' cohomology sequence, in general non-trivial in negative dimensions.

Let $\unsure{\overset{\infty}{M}|\widetilde M}=\operatorname{Hom}_{\mathbb Z}(M,\overline A^*)$ denote the Cartier dual of $M$. Here $\overline A^*$ denotes the group of invertible elements of $\overline A$. The $m$-th roots of $1$ are present so we do indeed have a perfect duality: $M\xrightarrow{\sim}\overset{\infty}{\overset{\infty}{M}}$, $\operatorname{Card} M=\operatorname{Card}\overset{\infty}{M}$, etc.

Let $\widehat H$ denote the \unsure{Pontrjagin|Pontryagin} dual of a locally compact group $H$.

\S\,2 \uline{Local Results}: Let $k$ be a locally compact field, i.e. $k=\mathbb R$, $k=\unsure{C|\mathbb C}$, or $k$ complete with respect to a discrete valuation, with finite residue class field. By local class field theory we have a canonical map $\operatorname{inv}_k:H^2(k,\overline k^*)\longrightarrow\mathbb Q/\mathbb Z$. Hence the cup product with respect to the canonical pairing \deleted{\illegible{1}} $M\times\overset{\infty}{M}\longrightarrow\overline k^*$ gives pairings
\[
H^r(k,M)\times H^{2-r}(k,\overset{\infty}{M})\longrightarrow\mathbb Q/\mathbb Z.
\]
These turn out to be perfect dualities of finite groups, i.e. we have

\uline{Theorem 2.1}: \uline{For all} $r\in\mathbb Z$, $H^r(k,M)$ \uline{is finite, and we have}
$H^{2-r}(k,M)\xrightarrow{\sim}\widehat{H^r(k,\overset{\infty}{M})}$.

If $k=\mathbb R$ or $C$ then $G_k$ is cyclic of order $2$ or $1$ and consequently $\operatorname{Card} H^r(k,M)$ is a power of $2$ and is independent of $r$.

\uline{Theorem 2.2}: \uline{If} $k$ \uline{is non-archimedean, then}
\[
\chi(k,M)\overset{\mathrm{defn.}}{=}
\frac{\operatorname{Card} H^0(k,M)\cdot\operatorname{Card} H^2(k,M)}{\operatorname{Card} H^1(k,M)}
=\frac{1}{(A:mA)}=|m|_k,
\]
\uline{where} $A$ \uline{is the valuation ring in} $k$, \uline{and where} $|m|_k$ \uline{denotes the ``normed'' abs. value of} $m$ \uline{in} $k$.
